\begin{ozetenv}
    Büyük Dil Modellerinin (BDM) ``Servis Olarak Makine Öğrenmesi'' (MLaaS) platformları aracılığıyla dağıtımı, kullanıcıların hassas bilgilerini güvenilmeyen sunuculara iletmesini gerektirdiğinden ciddi veri gizliliği endişeleri doğurmaktadır. Tam Homomorfik Şifreleme (THS), hesaplamaların doğrudan şifrelenmiş veri üzerinde yapılmasına olanak tanımakta, ancak Transformer mimarilerinin şifreli alana uyarlanması tüm polinom olmayan işlemlerin düşük dereceli polinom yaklaşımlarıyla değiştirilmesini gerektirmektedir.

    Bu tez, BERT'teki üç doğrusal olmayan işlemin---GELU, Softmax ve LayerNorm---tamamını öğrenilebilir Chebyshev polinom yaklaşımlarıyla değiştiren üç aşamalı aşamalı bir çerçeve olan Öğrenilebilir Polinom Yaklaşım Ağı'nı (LPAN) sunmaktadır. Temel katkılar şunlardır: (1)~standart Chebyshev'den $3,4\times$ daha düşük hata sağlayan dağılım ağırlıklı minimaks polinom başlangıç değerleri; (2)~dikkat düzeyinde bilgi damıtma ile aşamalı katman katman değiştirme; (3)~kafa uzmanlaşmasını mümkün kılan kafa başına polinom Softmax; ve (4)~ortak uyarlamalı ilk tam polinom LayerNorm değiştirme. SST-2 üzerinde dört BERT varyantında (Tiny'den Base'e) değerlendirilen LPAN, üç doğrusal olmayan işlemin tamamını değiştirirken yalnızca \%0,11--\%0,34 doğruluk düşüşü elde etmekte---yalnızca iki işlemi \%1--5 doğruluk düşüşüyle değiştiren mevcut yöntemlerden (THE-X, Iron, MPCFormer, BOLT) üstün performans göstermektedir. LPAN, tamamen etkileşimsiz şifreli Transformer çıkarımını mümkün kılan ilk çerçevedir.
\end{ozetenv}